% ===========================================================================
\section{The Correlated Altruist Bloc}
\label{sec:bayesian-bloc}

The fixed-point analyses of Sections~\ref{sec:equilibria}
and~\ref{sec:bayesian} model every author the same way: as an
individual best-responder, choosing one submission against a given
belief $\mathbf{f}$.
For own-payoff authors this is the whole story.
For total-welfare authors it undersells the objective: total welfare
is a property of the field, identical for every strategy in a given
tournament, so for an author choosing in isolation the rounds among
\emph{other} strategies are an additive constant --- which is why
they cancel in Equation~\eqref{eq:type-utility}, leaving even the
$\lambda = 1$ type optimizing only the rounds its own submission
plays.

Functional decision theory suggests a different accounting
\citep{yudkowsky2017functional}.
Total-welfare authors face the same decision problem with the same
evidence --- the same pool, the same matrices, the same objective ---
so their conclusions are logically correlated: whatever reasoning
carries one such author to a submission carries the others to it as
well.
An FDT author therefore does not ask ``what should I submit, taking
the other altruists' choices as given?''\ but ``what should authors
like me submit?'', treating the type's entire component of
$\mathbf{f}$ as a single decision variable.
This is the same correlation that strategies exploit \emph{inside}
the game --- a copy of my decision procedure decides as I do
(Sections~\ref{sec:properties} and~\ref{sec:related}) --- applied
one level up, to the choice of strategy itself.

A bloc of mass $m$ choosing jointly internalizes its influence on
the whole field, including the equilibrium responses of the other
types, and its objective is the field welfare
\begin{equation}
  V(\mathbf{f}) = \tfrac{1}{2}\, \mathbf{f}^\top W \mathbf{f},
  \label{eq:field-welfare}
\end{equation}
maximized over the bloc's component subject to
Equation~\eqref{eq:mixture-fixed-point} holding for the remaining
types.
Operationally: pin the bloc's mass on a candidate strategy, let the
other types re-equilibrate, and rank candidates by the resulting
$V(\mathbf{f}^*)$.

This is where the total-welfare objective departs from the
rounds-involving-me objective: the bloc's choice affects rounds it
never plays, through the equilibrium composition of the rest of the
field.
In particular, a bloc pinned on a strategy that punishes defection
suppresses the equilibrium mass that own-payoff types place on
defecting strategies, raising welfare in rounds among those types ---
a deterrence effect that is invisible to any per-strategy score.

\paragraph{The bloc avoids unconditional generosity.}
The bloc scan answers the question posed in
Section~\ref{sec:author-types}: candidate (a) loses to all three
conditional candidates.
Pinning the total-welfare bloc on \textbf{Generous} ranks near the
bottom of the pool: the own-payoff types' equilibrium response
includes strategies that exploit it, destroying welfare, while the
pin exerts no deterrent pressure.
The bloc's best choices are the conditional strategies; at large
$\beta$, \textbf{Cooperation coalition}, \textbf{Chivalry}, and
\textbf{Mirror} all attain the first-best $V = 2\ln(1 + E/2)$, and at
moderate $\beta$ the enforcing strategies win by exactly the
deterrence channel just described --- among the intuitive cases for
altruism, it is the \emph{enforcers'} that survives equilibrium
scrutiny.
The ranking is robust to the bloc's size
(Section~\ref{sec:sensitivity-types}).

\paragraph{Two decision theories, one conclusion.}
The bloc analysis is not a refinement of the fixed-point analysis
but an alternative to it.
The two rest on different decision theories at the author level ---
individual best response against a field taken as given, versus
correlated choice of an entire type --- and use different machinery
--- fixed-point iteration versus constrained optimization.
They nonetheless agree.
In the mixture fixed point, the total-welfare type's mass settles on
the conditionally cooperative core alongside the other types
(Section~\ref{sec:bayesian-illustration}); in the bloc scan, the
same conditional strategies top the ranking and unconditional
generosity finishes near the bottom.
The agreement is itself a finding: the conclusion that altruism is
best implemented conditionally does not depend on which decision
theory the authors use.
